\section{Falsifiable prediction QH1a}
\label{sec:qh1a_prediction}

Sections~\ref{sec:field_connection}--\ref{sec:crt_decoupling}
establish PT's structural position with respect to the
foundational debate: the holonomy angle is primitive, classical
objects are derivative, correlations between modular channels are
arithmetic and invariant under the emergent geometry. This
position is not merely interpretive: it entails a precise
experimental prediction, falsifiable on several existing
platforms, and structurally distinguishable from any prediction
of standard physics on a curved manifold.

\subsection{Statement of the prediction}
\label{ssec:qh1a_statement_en}

\begin{prediction}[QH1a: geometric invariance of MI]
\label{pred:qh1a_en}
In a multi-electronic atomic system probed by CRT spectroscopy
(channels $p = 3, 5, 7$), the mutual information
$I(n^{(p)}\,;\,n^{(q)})$ between modular occupation numbers:
\begin{enumerate}[label=(\roman*),nosep]
  \item is non-zero and given by the closed-form expression of
  Corollary~\ref{cor:closed_form_en} (arithmetic test);
  \item is invariant under any geometric transformation of the
  system (nuclear displacement, lattice deformation, detector
  separation) (geometric test).
\end{enumerate}
\end{prediction}

Condition (i) is the quantitative content of
Corollary~\ref{cor:closed_form_en}: the MI takes the values
$\mathcal{I}_{3,5} = 0.138$ bits and
$\mathcal{I}_{3,7} = 0.283$ bits (corpus values), predictable to
$\pm 20\%$ by the closed form with a single integer parameter
$k_{p,q}$ per pair. Any measurement failing to reproduce these
values (within the statistical precision of the experiment)
invalidates the prediction.

Condition (ii) is the structural content of
Theorem~\ref{thm:invariance_geom_en}: the MI must remain constant
under any reparametrisation of the spatial geometry. Any
measurement detecting a dependence of the MI on detector
separation, internuclear distance, or any other geometric
parameter, invalidates the prediction.

\subsection{Falsification criterion}
\label{ssec:qh1a_falsification_en}

Let $G_0, G_1$ be two distinct geometric configurations of the
same quantum system, and let
$\mathcal{I}_{p,q}(G_0), \mathcal{I}_{p,q}(G_1)$ be the measured
mutual information under each. Define
\begin{equation}
  \Delta\mathcal{I}_{p,q}
  \;:=\;
  \mathcal{I}_{p,q}(G_1) - \mathcal{I}_{p,q}(G_0)
  \label{eq:Delta_I_def_en}
\end{equation}
and let $\epsilon_{\rm stat}$ be the statistical uncertainty of
$\Delta\mathcal{I}_{p,q}$ as evaluated by the measurement
procedure. PT predicts
\begin{equation}
  \bigl|\Delta\mathcal{I}_{p,q}\bigr| \;<\; \epsilon_{\rm stat}
  \qquad \text{for every pair }(G_0, G_1).
  \label{eq:QH1a_predict_en}
\end{equation}
\emph{Falsification criterion.} If an experiment measures
$|\Delta\mathcal{I}_{p,q}| > 5 \epsilon_{\rm stat}$ for a pair of
distinct geometric configurations, prediction QH1a is falsified,
and with it the invariance
theorem~\ref{thm:invariance_geom_en}. Such a falsification would
also be a direct refutation of PT, since $\theta_p$ and
$\mathcal{A}_p$ are structural objects of the theory and cannot
be adjusted to absorb a geometric signal.

\subsection{Candidate experimental platforms}
\label{ssec:qh1a_platforms_en}

Three contemporary platforms offer access to the modular
occupation numbers $n^{(p)}$ with independent control of the
spatial geometry.

\paragraph{(a) Cold atoms in optical lattices.}
An optical lattice with variable spacing allows one to scan the
distance between adjacent potential wells while measuring the
modular occupation distribution of internal levels by fluorescence
spectroscopy. Geometry is controlled by the wavelength of the
trapping laser; the occupation numbers $n^{(p)}$ are accessible
via resonant measurement at the modular residue level. Expected
precision: $\epsilon_{\rm stat} \sim 10^{-3}$ bits on
$\mathcal{I}_{3,5}$ for $N \sim 10^4$ trapped atoms, ample
sensitivity to falsify QH1a at $5\sigma$ in the presence of
geometric coupling.

\paragraph{(b) Trapped ions at controlled separations.}
A chain of Yb$^+$ or Be$^+$ ions allows one to vary the inter-ion
separation with micrometre precision while measuring inter-channel
correlations by state tomography. The hyperfine structure of
internal levels provides the modular occupation numbers for
$p = 3$ (effective SU(3) colour channel) and $p = 5$ (flavour
channel in configurations with several valence electrons).
Expected precision: $\epsilon_{\rm stat} \sim 10^{-4}$ bits, the
tightest constraint of the three.

\paragraph{(c) Diatomic molecules at adjustable internuclear distance.}
A diatomic molecule probed by laser spectroscopy at sub-MHz
resolution allows one to vary the internuclear distance $R$ via
an external electric field (Stark effect) and to measure
simultaneously the modular occupation numbers of the electronic
levels. The signature sought is the absence of dependence of
$\mathcal{I}_{3,5}$ on $R$ over a range of several angstroms.
Expected precision: $\epsilon_{\rm stat} \sim 5 \times 10^{-3}$
bits, more modest but structurally complementary (the geometric
coupling expected of a standard theory is here dominated by the
vibrational-potential variation).

\subsection{Contrast with standard physics}
\label{ssec:qh1a_contrast_en}

A standard quantum theory on curved manifold predicts generically
$|d\mathcal{I}_{p,q}/dG| \neq 0$. The structural reason is that
matter fields are, in the standard formalism, sections of bundles
above the curved base, and their correlations inherit the metric
structure of the base via minimal gravitational couplings. The
effect is typically weak (suppressed by $1/M_{\rm Pl}^2$ or by
the local gravitational coupling), but it is \emph{non-zero} and
computable.

PT, by contrast, predicts $d\mathcal{I}_{p,q}/dG = 0$ identically.
The reason is that the modular channels $\mathcal{A}_p$ are
defined before the metric $G$ and are measurably anterior to it.
No function of the metric appears in the closed
form~\eqref{eq:closed_form_en}; the derivative with respect to $G$
is exactly zero.

This contrast makes QH1a a strong discriminating test between the
two frameworks. A positive measurement ($|d\mathcal{I}/dG| >
5\epsilon$) falsifies PT and confirms the standard framework; a
null measurement ($|d\mathcal{I}/dG| < \epsilon$) confirms PT at
the structural level and excludes any minimal geometric coupling
above the achieved precision. None of the current platforms has
yet performed the complete measurement --- in particular, none
has scanned the geometric dependence of the MI over several
orders of magnitude of separation. This is the principal
experimental gap to be closed in order to subject PT to a
genuine falsificationist test.
